\clearpage
\section*{LAMPIRAN A Template TOR}

\begin{framed}
\begin{center}
{\bfseries\Large TERM OF REFERENCES}
\end{center}

\textit{Term of Reference} (TOR) merupakan dokumen yang berisi pernyataan kebutuhan pengguna (\textit{user}) yang sangat beragam yaitu mulai dari layer yang paling umum sampai dengan layer yang detail. Untuk Kerja Praktik (KP), TOR umumnya dibuat oleh instansi/perusahaan pemberi pekerjaan. Hal ini untuk menegaskan apa yang dibutuhkan perusahaan, apa yang harus dikerjakan oleh mahasiswa KP untuk mewujudkan kebutuhan perusahaan. TOR antara lain berisi sub bab sebagai berikut:

\begin{enumerate}[label=\arabic*.,leftmargin=1.2cm]
\item \textbf{TUJUAN PEKERJAAN}

Tujuan Pekerjaan merupakan pernyataan objektif pekerjaan yang dapat berupa hasil/produk, pekerjaan yang harus dilaksanakan, atau dampak pekerjaan dan kepada siapa hasil pekerjaan tersebut diperuntukan. Sub bab ini bisa juga diberi judul Latar Belakang.

\item \textbf{LINGKUP PEKERJAAN}

Lingkup Pekerjaan merupakan pernyataan tentang cakupan pekerjaan yang harus dilakukan untuk mencapai tujuan. Lingkup di sini mungkin lebih besar dari lingkup pekerjaan yang harus dilakukan oleh mahasiswa yang sedang melaksanakan KP.

\item \textbf{HASIL PEKERJAAN}

Hasil Pekerjaan merupakan pernyataan deskripsi produk yang dihasilkan atau dampak pekerjaan. Contoh dampak pekerjaan adalah SDM yang andal.

\item \textbf{BATASAN}

Batasan-batasan merupakan pernyataan-pernyataan yang bersifat membatasi lingkup atau hasil untuk mengamankan kedua belah pihak baik di pihak pengembang maupun klien, dapat berupa batasan terhadap proses bisnis, hasil, lingkungan pengembangan, lingkungan operasi/sistem, dll. Batasan juga berisi lingkup pekerjaan yang harus dilakukan oleh mahasiswa KP. Jika lingkup pekerjaan memang harus dilakukan semua oleh mahasiswa KP, maka batasan ini berisi spesifikasi lain yang dapat membantu mahasiswa kerja praktek dalam menyelesaikan pekerjaannya. Batasan lingkungan operasi/sistem adalah spesifikasi terhadap perangkat yang akan digunakan.

\textit{Disarankan ditambahkan gambar untuk lebih menunjukkan koordinat/posisi hasil dan/atau pekerjaan dengan keterkaitan terhadap sistem global (tergantung project) atau interaksi antar modul satu layer di bawahnya. Gambar ini dapat diletakkan di bagian tujuan, atau lingkup atau hasil atau ditambahkan di bagian terpisah.}

\item \textbf{ASUMSI}

Asumsi-asumsi merupakan pernyataan tentang status dan kondisi yang disepakati di awal untuk pelaksanaan pekerjaan, termasuk di dalamnya kebergantungan dan kendala.

\item \textbf{METODOLOGI PEKERJAAN}

Metodologi adalah strategi dan tahapan untuk melaksanakan pekerjaan yang mencakup cara atau tahapan, alat dan hasil antara setiap tahapan atau pernyataan yang mendeskripsikan bagaimana lingkup dan hasil bisa dicapai. Metodologi pekerjaan mencakup:
\begin{enumerate}[label=\roman*.,leftmargin=1.2cm]
\item penjelasan singkat mengenai metodologi dan
\item pemetaan metodologi ke dalam tahap-tahap yang digunakan.
\end{enumerate}

\item \textbf{LINGKUNGAN PENGEMBANGAN}

Lingkungan pengembangan dapat meliputi spesifikasi perangkat keras dan perangkat lunak yang digunakan untuk mengembangkan sistem yang akan dibuat di dalam KP.

\item \textbf{JADWAL}

Jadwal merupakan pernyataan rencana waktu dalam pelaksanaan pekerjaan yang dirinci untuk setiap tahapan pekerjaan. Jadwal ini dituangkan dalam bentuk GANTT chart.

\item \textbf{PENUGASAN STAF}

Penugasan Staf merupakan pernyataan pelimpahan wewenang kepada staf yang ditunjuk untuk melaksanakan pekerjaan sesuai dengan jabatan masing-masing. Bagian ini terutama dibutuhkan oleh mahasiswa yang mengerjakan KP berkelompok dan tiap-tiap mahasiswa memiliki tanggung jawab untuk tahapan atau bagian pekerjaan tertentu (walaupun masih dalam satu topik).

\item \textbf{MEKANISME PENGENDALIAN DAN PEMANTAUAN}

Mekanisme pemantauan dan pengendalian merupakan pernyataan tata cara pemantauan dan pengendalian dalam pelaksanaan pekerjaan. Mekanisme diutamakan mengenai mekanisme pelaporan, meliputi laporan kemajuan proyek dan frekuensi pelaporan, serta alur pelaporan tersebut (dari siapa ke siapa). Bagian ini akan banyak berhubungan dengan pencatatan aktivitas kegiatan (\textit{log activity}).

\textit{Ditambahkan milestone untuk menyatakan penyetoran laporan, presentasi, demo.}

\item \textbf{MANAJEMEN RESIKO (OPTIONAL)}

Manajemen resiko merupakan pernyataan terhadap komitmen kita terhadap pekerjaan, yang dimulai dengan identifikasi resiko yang ada dan cara pengendalian resiko tersebut. Jika memungkinkan menggunakan pembobotan untuk menetapkan prioritas.

\item \textbf{LAIN-LAIN}

Jika ada hal lain yang perlu disetujui antara kedua belah pihak (perusahaan dan mahasiswa kerja praktik), maka informasi tersebut dapat dituliskan di bagian ini. Salah satu contoh adalah mengenai insentif yang akan diberikan pada mahasiswa KP.
\end{enumerate}

\vspace{0.8cm}
\begin{center}
Disetujui oleh,

\vspace{1.2cm}
<Pembimbing Instansi/Perusahaan>

\vspace{1.2cm}
Mahasiswa Kerja Praktik

\vspace{1.5cm}
\begin{tabular}{ccc}
<Mahasiswa 1> & \hspace{2cm} <Mahasiswa 2> & \hspace{2cm} <Mahasiswa 3>
\end{tabular}
\end{center}
\end{framed}

\clearpage
\chapter*{Lampiran B. <Lampiran Lain>}
\addcontentsline{toc}{chapter}{Lampiran B. <Lampiran Lain>>}

\textit{Lampiran B bersifat optional untuk mahasiswa yang melakukan KP. Silakan tambahkan lampiran lain di bagian ini.}

\clearpage